% Gemini theme
% https://github.com/anishathalye/gemini

\documentclass[final]{beamer}

% ====================
% Packages
% ====================

\usepackage[T1]{fontenc}
\usepackage{lmodern}

% CAMBIO 1: Tamaño A0 y orientación portrait (vertical)
\usepackage[orientation=portrait,size=a0,scale=1.05]{beamerposter}

\usetheme{gemini}
\usecolortheme{argentina}
\usepackage{graphicx}
\usepackage{booktabs}
\usepackage{tikz}
\usepackage{pgfplots}
\pgfplotsset{compat=1.14}
\usepackage{anyfontsize}
\usepackage{pgfplots}
\pgfplotsset{compat=1.18}
\usepackage{xcolor}

% Colores aproximados a los del gráfico original
\definecolor{graymid}{RGB}{128,128,128}       % gris para 1a
\definecolor{tealmid}{RGB}{0,128,128}         % verde azulado medio
\definecolor{tealdark}{RGB}{0,96,96}          % verde azulado oscuro
\definecolor{coralmid}{RGB}{255,127,80}       % coral para 2a y línea punteada
\definecolor{coraldark}{RGB}{205,92,92}       % coral oscuro para texto

% ====================
% Lengths
% ====================

% CAMBIO 2: Matemáticas para 3 columnas
\newlength{\sepwidth}
\newlength{\colwidth}
\setlength{\sepwidth}{0.025\paperwidth}
\setlength{\colwidth}{0.3\paperwidth}
\newlength{\widecolwidth}
\setlength{\widecolwidth}{2\colwidth + \sepwidth}

\newenvironment{labelblock}[2]{%
  \begin{block}{%
      {\small\color{structure.fg!70!bg}#1}%   etiqueta pequeña (grisácea)
      \\[2pt]                                  % salto con pequeño espacio
      {\large\bfseries#2}%                     % título grande y negrita
      }%
      }{%
  \end{block}%
  }

  \newcommand{\separatorcolumn}{\begin{column}{\sepwidth}\end{column}}

    % ====================
    % Title
    % ====================

    \title{Resilient food solutions for Argentina under volcanic abrupt sunlight reduction scenarios: a linear optimization approach}

    \author{Ulloa Ruiz, M.A. · Torres Celis, J.A. · Rivers, M. · Denkenberger, D.C. · García Martínez, J.B.}

    \institute[shortinst]{Observatorio de Riesgos Catastróficos Globales (ORCG) · ALLFED · University of Canterbury}

    % ====================
    % Footer (optional)
    % ====================

    \footercontent{
      Ulloa Ruiz et al. · ORCG / ALLFED · 2024 · REDER 8(2) \hfill
      Volcanic Impacts on Climate and Society Conference 2026}

    % ====================
    % Body
    % ====================
    \begin{document}

    \begin{frame}[t]
      % ==================== PRIMERA FILA (tres columnas normales) ====================
      \begin{columns}[t]
        \separatorcolumn

        \begin{column}{\colwidth}
          \begin{labelblock}{Background}{Motivation}
            Large-scale volcanic eruptions inject sulfate aerosols into the stratosphere, reducing incoming solar radiation for months to years — an Abrupt Sunlight Reduction Scenario (ASRS). The 1815 Tambora eruption and subsequent "year without a summer" offers the clearest historical precedent of the resulting agricultural collapse and food crises across the Northern Hemisphere.

            Modern climate models confirm that eruptions generating $5-150$ Tg of stratospheric particulate can reduce global caloric crop production by $7-90\%$, depending on magnitude. Because existing food stocks can sustain populations for only a few months, a rapid adaptive response is required. This study focuses on Argentina as a strategic case study for the Southern Hemisphere.
          \end{labelblock}
        \end{column}

        \separatorcolumn

        \begin{column}{\colwidth}
          \begin{labelblock}{Research aim}{Objective}
            To evaluate the effectiveness of a phased portfolio of adaptive food-system interventions on per-capita caloric availability in Argentina across a seven-year post-eruption window, modeled under volcanic ASRS conditions spanning $5$ to $150$ Tg of atmospheric particulate loading.

            Specific aims:
            \begin{enumerate}
              \item Quantify production loss under volcanic ASRS without adaptation
              \item Model phased interventions and their combined caloric yield
              \item Assess regional food-security and export capacity implications
            \end{enumerate}
          \end{labelblock}
        \end{column}

        \separatorcolumn

        \begin{column}{\colwidth}
          \begin{labelblock}{At a glance}{Key results}
            \definecolor{coral}{RGB}{216,90,48}   % #D85A30
            \definecolor{amber}{RGB}{186,117,23}  % #BA7517
            \definecolor{teal}{RGB}{29,158,117}   % #1D9E75

            \noindent
            \begin{minipage}[t]{0.48\linewidth}
              \centering
              {\sffamily\fontsize{22}{26}\selectfont\bfseries\textcolor{coral}{70\%}}\par
              \vspace{2pt}
              {\sffamily\footnotesize\color{black!70}Population fed without adaptation (150 Tg)}
            \end{minipage}\hfill
            \begin{minipage}[t]{0.48\linewidth}
              \centering
              {\sffamily\fontsize{22}{26}\selectfont\bfseries\textcolor{teal}{3-6$\times$}}\par
              \vspace{2pt}
              {\sffamily\footnotesize\color{black!70}Min. requirement met with full adaptation}
            \end{minipage}

            \vspace{6pt}

            \noindent
            \begin{minipage}[t]{0.48\linewidth}
              \centering
              {\sffamily\fontsize{22}{26}\selectfont\bfseries\textcolor{amber}{300M}}\par
              \vspace{2pt}
              {\sffamily\footnotesize\color{black!70}People in the region who could receive surplus exports}
            \end{minipage}\hfill
            \begin{minipage}[t]{0.48\linewidth}
              \centering
              {\sffamily\fontsize{22}{26}\selectfont\bfseries\textcolor{teal}{7 yr}}\par
              \vspace{2pt}
              {\sffamily\footnotesize\color{black!70}Modeling window post-event}
            \end{minipage}

            \vspace{6pt}
            \small
            All scenarios modeled against a 2,100 kcal/day/person minimum requirement. Under the most severe scenario (150 Tg), full intervention brings net production to 7,800-14,000 kcal/day/person.
          \end{labelblock}
        \end{column}

        \separatorcolumn
      \end{columns}

      % Espacio vertical entre la primera y la segunda fila
      \vspace{0.5cm}

      % ==================== SEGUNDA FILA (bloque ancho que ocupa 2 columnas + 1 columna independiente) ====================
      \begin{columns}[t]
        \separatorcolumn

        \begin{column}{\widecolwidth}
          \begin{labelblock}{Methods}{Methodology}
            % Contenido para el bloque Methods – Methodology
            \begin{minipage}[t]{0.48\linewidth}
              \raggedright
              \vspace{0pt} % para alineación superior precisa
              {\small\bfseries Model framework}

              \small
              Adapted from the ALLFED Integrated Food System Model (Rivers et al., 2023),
              implemented in Python. Linear optimization minimizes starvation while
              maximizing caloric availability across the population, distributing food
              stocks and new production monthly over 120 months.

              \vspace{3pt}
              Crop yield loss estimates drawn from Xia et al.\ (2022), covering six
              particulate loading levels (5--150\,Tg). Baseline production data from
              FAO (2020--2022).
            \end{minipage}\hfill
            \begin{minipage}[t]{0.48\linewidth}
              \raggedright
              {\small\bfseries Scenarios modeled}

              \vspace{4pt}
              {\small
              \begin{tabular}{@{} l p{\dimexpr\linewidth-6.5cm\relax} c @{}}
                \toprule
                Scenario & Description & Adapt. \\
                \midrule
                1a / 1b & Pre-event baseline (net / gross) & \textcolor{red!70!black}{\textsf{none}} \\
                2a      & 150\,Tg, no adaptation               & \textcolor{red!70!black}{\textsf{none}} \\
                2b      & Redirection from livestock \& biofuels & \textcolor{orange}{\textsf{partial}} \\
                2c      & + resilient food solutions            & \textcolor{blue}{\textsf{moderate}} \\
                2d      & + greenhouse scale-up \& area expansion & \textcolor{green!50!black}{\textsf{full}} \\
                \bottomrule
              \end{tabular}
              }

              \vspace{6pt}
              \small
              Country selection via multi-criteria decision matrix: supply-chain resilience
              (WEF GCI 2020), governance effectiveness (World Bank WGI 2021), and ASRS food
              production capacity.
            \end{minipage}
          \end{labelblock}

        \end{column}

        \separatorcolumn

        % Tercera columna (ancho normal)

        \begin{column}{\colwidth}
          \begin{labelblock}{Case selection}{Why Argentina?}
            \definecolor{tealdark}{HTML}{006666}
            \begin{minipage}[t]{0.48\linewidth}
              \raggedright
              {\footnotesize\bfseries\textcolor{tealdark}{Southern hemisphere location}}\\
              {\scriptsize Temperate zone; less exposed to volcanic winter cooling concentrated in the
              Northern Hemisphere, as Tambora (1815) demonstrated.}
              \vspace{5pt}

              {\footnotesize\bfseries\textcolor{tealdark}{Diversified economy}}\\
              {\scriptsize Strong agri-industrial base; capacity to repurpose paper mills, biorefineries,
              and livestock infrastructure rapidly.}
            \end{minipage}\hfill
            \begin{minipage}[t]{0.48\linewidth}
              \raggedright
              {\footnotesize\bfseries\textcolor{tealdark}{Agricultural capacity}}\\
              {\scriptsize 154M\,ha total agricultural land; 43M\,ha easily cultivable. Pre-ASRS gross
              production {\textasciitilde}28,000\,kcal/day/person.}
              \vspace{5pt}

              {\footnotesize\bfseries\textcolor{tealdark}{Regional connectivity}}\\
              {\scriptsize Rail links with 6 neighbours; maritime trade with Brazil (key fossil fuel and
              food exporter), enabling surplus distribution.}
            \end{minipage}

            \vspace{6pt}
            {\small Uruguay and Chile were also identified as priority candidates in the regional
            prioritization matrix. Argentina selected for its higher production capacity under modeled
            scenarios.}
          \end{labelblock}
        \end{column}

        \separatorcolumn
      \end{columns}

      \vspace{0.5cm}

      % ==================== TERCERA FILA (ocupa las 3 columnas completas) ====================
      \begin{columns}[t]
        \separatorcolumn
        % Columna que abarca el ancho total de las tres columnas más los dos separadores intermedios
        \newlength{\fullwidth}
        \setlength{\fullwidth}{3\colwidth + 2\sepwidth}
        \begin{column}{\fullwidth}
          \begin{labelblock}{Adaptive strategies}{Interventions evaluated}
            % Colores de las barras izquierdas
            \definecolor{tealmid}{HTML}{2A9D8F}
            \definecolor{ambermid}{HTML}{DBA100}

            % ── Primera fila (4 columnas) ──
            \begin{minipage}[t]{0.22\linewidth}
              \raggedright
              \textcolor{tealmid}{\vrule width 3pt}\hspace{3pt}%
              \begin{minipage}[t]{\dimexpr\linewidth-6pt\relax}
                {\bfseries\small Livestock \& biofuel redirection}\par\vspace{2pt}
                {\footnotesize Scaling down cattle and biofuel production releases grain
                currently consumed by animals. Estimated release: up to
                {\textasciitilde}15,000\,kcal/capita/day. Single highest-impact
                intervention (+300\% availability vs.\ scenario 2a).}
              \end{minipage}
            \end{minipage}\hfill
            \begin{minipage}[t]{0.22\linewidth}
              \raggedright
              \textcolor{tealmid}{\vrule width 3pt}\hspace{3pt}%
              \begin{minipage}[t]{\dimexpr\linewidth-6pt\relax}
                {\bfseries\small Cold-tolerant crop relocation}\par\vspace{2pt}
                {\footnotesize Introducing northern-hemisphere cold-tolerant varieties
                (wheat, potato, oat, rapeseed) to latitudes where soy and maize yields
                collapse. Relocation should begin immediately after event.}
              \end{minipage}
            \end{minipage}\hfill
            \begin{minipage}[t]{0.22\linewidth}
              \raggedright
              \textcolor{ambermid}{\vrule width 3pt}\hspace{3pt}%
              \begin{minipage}[t]{\dimexpr\linewidth-6pt\relax}
                {\bfseries\small Low-tech greenhouse scale-up}\par\vspace{2pt}
                {\footnotesize Rapid expansion from 6,500\,ha to {\textasciitilde}2.3M\,ha
                within 2 years ($\times$360). Controls temperature variability under
                reduced insolation. Scenario 2d key driver.}
              \end{minipage}
            \end{minipage}\hfill
            \begin{minipage}[t]{0.22\linewidth}
              \raggedright
              \textcolor{ambermid}{\vrule width 3pt}\hspace{3pt}%
              \begin{minipage}[t]{\dimexpr\linewidth-6pt\relax}
                {\bfseries\small Cultivated area expansion}\par\vspace{2pt}
                {\footnotesize From 31M\,ha (current 5 crops) to 43M\,ha (easily
                cultivable) or up to 108M\,ha. With expansion to 108M\,ha, gross
                production could sustain 6--22$\times$ national population.}
              \end{minipage}
            \end{minipage}

            \vspace{6pt}

            % ── Segunda fila (4 columnas) ──
            \begin{minipage}[t]{0.22\linewidth}
              \raggedright
              \textcolor{tealmid}{\vrule width 3pt}\hspace{3pt}%
              \begin{minipage}[t]{\dimexpr\linewidth-6pt\relax}
                {\bfseries\small Seaweed cultivation revival}\par\vspace{2pt}
                {\footnotesize Argentine seaweed industry nearly extinct since mid-20th
                century. High potential along Chubut coast. Knowledge transfer from
                countries with active algae industries required.}
              \end{minipage}
            \end{minipage}\hfill
            \begin{minipage}[t]{0.22\linewidth}
              \raggedright
              \textcolor{tealmid}{\vrule width 3pt}\hspace{3pt}%
              \begin{minipage}[t]{\dimexpr\linewidth-6pt\relax}
                {\bfseries\small Rationing \& waste reduction}\par\vspace{2pt}
                {\footnotesize Assumed in baseline model: equitable distribution across
                population; reduction of food loss due to scarcity. Necessary for all
                scenarios to prevent unequal access collapse.}
              \end{minipage}
            \end{minipage}\hfill
            \begin{minipage}[t]{0.22\linewidth}
              \raggedright
              \textcolor{ambermid}{\vrule width 3pt}\hspace{3pt}%
              \begin{minipage}[t]{\dimexpr\linewidth-6pt\relax}
                {\bfseries\small Industrial sugar (lignocellulosic)}\par\vspace{2pt}
                {\footnotesize Rapid repurposing of pulp and paper mills into
                lignocellulosic sugar factories. Contributes
                {\textasciitilde}200\,kcal/person/day in months 15--35
                (scenario 2c/2d).}
              \end{minipage}
            \end{minipage}\hfill
            \begin{minipage}[t]{0.22\linewidth}
              \raggedright
              \textcolor{ambermid}{\vrule width 3pt}\hspace{3pt}%
              \begin{minipage}[t]{\dimexpr\linewidth-6pt\relax}
                {\bfseries\small Single-cell protein (SCP)}\par\vspace{2pt}
                {\footnotesize Methane-based SCP production to compensate animal protein
                losses. Used as animal feed substitute in scenarios 2c/2d to free crop
                production for human consumption.}
              \end{minipage}
            \end{minipage}

            \vspace{8pt}

            % ── Leyenda ──
            \begin{minipage}{\linewidth}
              \footnotesize
              \raisebox{1pt}{\textcolor{tealmid}{\rule{10pt}{10pt}}}
              \hspace{3pt} Supply redirection \& agricultural adaptation
              \hspace{24pt}
              \raisebox{1pt}{\textcolor{ambermid}{\rule{10pt}{10pt}}}
              \hspace{3pt} Industrial \& high-technology solutions
            \end{minipage}
          \end{labelblock}
        \end{column}
        \separatorcolumn
      \end{columns}

      \vspace{0.5cm}

      % ==================== TERCERA FILA (bloque ancho que ocupa 2 columnas + 1 columna independiente) ====================
      \begin{columns}[t]
        \separatorcolumn

        % Columna que ocupa el espacio de las dos primeras columnas originales + el separador entre ellas
        \begin{column}{\widecolwidth}
          \begin{labelblock}{Results}{Caloric availability by scenario}
            \begin{tikzpicture}
              \begin{axis}[
                  width=0.95\linewidth,
                  height=0.4\linewidth,
                  symbolic x coords={1a,1b,2a,2b,2c,2d},
                  xtick={1a,1b,2a,2b,2c,2d}, % <-- CORRECCIÓN 1: Declarar explícitamente las marcas
                  xticklabel style={font=\footnotesize},
                  ytick={0,5000,10000,20000,30000},
                  yticklabels={0,5k,10k,20k,30k},
                  ymin=0, ymax=31000,
                  axis lines=left,
                  enlarge x limits=0.15,
                  bar width=5cm, % <-- CORRECCIÓN 3: Ajustado a un tamaño que no desborde el documento
                  grid=major,
                  grid style={very thin, gray!40},
                  major grid style={gray!40},
                  tick label style={font=\footnotesize},
                  axis line style={thin, gray!60},
                  ylabel={kcal/day/person},
                  ylabel style={font=\footnotesize},
                ]

                % Barras existentes
                \addplot[ybar, fill=graymid, opacity=0.7, rounded corners=2pt] coordinates {(1a,4800)};
                \addplot[ybar, fill=tealmid, opacity=0.55, rounded corners=2pt] coordinates {(1b,28000)};
                \addplot[ybar, fill=coralmid, opacity=0.9, rounded corners=2pt] coordinates {(2a,1500)};
                \addplot[ybar, fill=tealmid, opacity=0.75, rounded corners=2pt] coordinates {(2b,7800)};
                \addplot[ybar, fill=tealdark, opacity=0.8, rounded corners=2pt] coordinates {(2c,10000)};
                \addplot[ybar, fill=tealdark, opacity=0.95, rounded corners=2pt] coordinates {(2d,14000)};

                % Línea de umbral corregida (2100 kcal)
                % <-- CORRECCIÓN 2: Usar \draw combinando coordenadas relativas (rel axis cs)
                % con las absolutas (axis cs) para forzar que ocupe todo el margen horizontal.
                \draw[dashed, thick, gray!70]
                ({rel axis cs:0,0} |- {axis cs:1a,2100}) --
                ({rel axis cs:1,0} |- {axis cs:1a,2100})
                node[midway, above, font=\footnotesize, text=gray!70] {2,100};

                % Etiquetas sobre las barras
                \node at (axis cs:1a,4800) [above, font=\footnotesize, black]      {4.8k};
                \node at (axis cs:1b,28000)[above, font=\footnotesize, tealdark]   {28k};
                \node at (axis cs:2a,1500) [above, font=\footnotesize, coraldark]  {1.5k};
                \node at (axis cs:2b,7800) [above, font=\footnotesize, tealdark]   {7.8k};
                \node at (axis cs:2c,10000)[above, font=\footnotesize, tealdark]   {10k};
                \node at (axis cs:2d,14000)[above, font=\footnotesize, tealdark]   {14k};

              \end{axis}
            \end{tikzpicture}
          \end{labelblock}
        \end{column}

        \separatorcolumn

        % Tercera columna (ancho normal)

        \begin{column}{\colwidth}
          \begin{labelblock}{Discussion}{Conclusions}
            \begin{itemize}
              \item A volcanic ASRS of $150$ Tg would reduce Argentina's net food production to $\~{} 70\%$ of minimum population requirements - famine without intervention, even with rationing alone.
              \item Redirecting food from livestock and biofuels - the most impactful single measure - can alone increase human-available calories by up to $300\%$, bringing production above the survival threshold.
              \item A phased combination of all modeled interventions ($2d$) produces $3-6\times$ minimum national requirements, enabling surplus export to support up to $300$ million additional people in the region.
              \item Southern Cone countries occupy a structurally advantageous position in any ASRS scenario, including volcanic origin; advance national planning is both technically feasible and strategically urgent.
              \item Success depends on pre-event institutional preparedness: contingency plans, strategic input stockpiles, legislative flexibility, and inter-provincial coordination within Argentina's federal structure.
            \end{itemize}
          \end{labelblock}
        \end{column}

        \separatorcolumn
      \end{columns}
      \vspace{0.5cm}

      % ==================== CUARTA FILA (tres columnas normales) ====================
      \begin{columns}[t]
        \separatorcolumn

        \begin{column}{\colwidth}
          \begin{labelblock}{Caveats}{Key uncertainties}
            Results are sensitive to assumptions that require further research:
            \begin{itemize}
              \item Volcanic aerosol optical depth
              \item Regional agroclimate response
              \item Algae scale-up feasibility
              \item Nutritional adequacy of resilient diets
              \item International trade continuity
              \item Federal coordination heterogeneity
              \item Demographic nutritional variation
              \item Industrial system disruption depth
            \end{itemize}

            Yield loss estimates approximate for cold-sensitive crops (maize, soy) vs. cold-tolerant crops (potato, wheat, oat). Local edaphological and hydrological variation across Argentine provinces not yet captured.
          \end{labelblock}
        \end{column}

        \separatorcolumn

        \begin{column}{\colwidth}
          \begin{labelblock}{Agenda}{Future work}
            Priority research directions to reduce uncertainty and increase policy relevance:
            \begin{itemize}
              \item Argentina-specific agroclimate models calibrated to volcanic forcing scenarios
              \item Province-level food distribution and economic modeling
              \item Pilot trials for low-tech greenhouse deployment and algae cultivation in Patagonia
              \item Subgroup-specific nutritional needs analysis under ASRS conditions
              \item Comparative analysis: volcanic vs. nuclear ASRS forcing on Southern Cone food systems
            \end{itemize}
          \end{labelblock}
        \end{column}

        \separatorcolumn

        \begin{column}{\colwidth}
          \begin{labelblock}{References \& data}{Selected references}
            \begin{thebibliography}{9}
              \bibitem{xia2022}
                Xia et al. (2022).
                \newblock Global food insecurity and famine from reduced crop production due to nuclear war soot.
                \newblock \emph{Nature Food}, 3, 586--596.

              \bibitem{rivers2022}
                Rivers et al. (2022).
                \newblock Food system adaptation and maintaining trade greatly mitigate global famine in ASRS.
                \newblock ALLFED Preprint.

              \bibitem{alvarado2020}
                Alvarado et al. (2020).
                \newblock Scaling of greenhouse crop production in low sunlight scenarios.
                \newblock \emph{Sci. Total Environ.}, 707.

              \bibitem{wilson2023}
                Wilson et al. (2023).
                \newblock Impact of the Tambora eruption of 1815 on islands.
                \newblock \emph{Sci. Rep.}, 13, 3649.

              \bibitem{ulloa2024}
                Ulloa Ruiz et al. (2024).
                \newblock Soluciones alimentarias resilientes...
                \newblock \emph{REDER}, 8(2), 159--176.
            \end{thebibliography}

            \vspace{0.5cm}
            \scriptsize
            Model code: \url{https://github.com/allfed/allfed-integrated-model} \\
            Supplementary material: \url{https://doi.org/10.5281/zenodo.11658966} \\
            Correspondence: \href{mailto:juan@allfed.info}{juan@allfed.info}
          \end{labelblock}
        \end{column}

        \separatorcolumn
      \end{columns}

    \end{frame}

  \end{document}
